\chapter{Systems of Linear Equations}

\chapterintro{A system of two linear equations asks where two lines meet, and the answer is a point, no point, or every point on a shared line. The elimination method you learn here is the $2 \times 2$ case of Gaussian elimination, which is the workhorse of Course~6.}

\section{Core Ideas}

\paragraph{What a solution is} A pair $(x, y)$ that satisfies both equations at once. Graphically it is the intersection of the two lines.

\paragraph{Substitution} Solve one equation for a variable, substitute that expression into the other equation, solve, then back-substitute. Best when one equation is already solved for a variable.

\paragraph{Elimination} Multiply one or both equations by constants so that one variable has opposite coefficients, then add the equations to eliminate it. These are row operations: scaling a row and adding one row to another never changes the solution set.

\paragraph{Three outcomes} Lines with different slopes cross once (one solution). Parallel distinct lines never cross (no solution; elimination produces a false statement like $0 = 5$). The same line written twice gives infinitely many solutions (elimination produces $0 = 0$).

\paragraph{The determinant test} For $a_1 x + b_1 y = c_1$ and $a_2 x + b_2 y = c_2$, the number $D = a_1 b_2 - a_2 b_1$ decides it. If $D$ is not zero there is exactly one solution, given by Cramer's rule.

\begin{keybox}{Cramer's rule for a $2 \times 2$ system}
\[
x = \frac{c_1 b_2 - c_2 b_1}{a_1 b_2 - a_2 b_1}, \qquad y = \frac{a_1 c_2 - a_2 c_1}{a_1 b_2 - a_2 b_1}
\]
\end{keybox}

\paragraph{Back substitution} A triangular system such as $x + y + z = 6$, $2y + z = 7$, $z = 3$ is solved from the bottom up. Elimination's whole purpose is to produce this shape.

\section{Worked Examples}

\begin{example}
Solve $y = 2x - 1$ and $3x + y = 9$ by substitution.
\solution Replace $y$: $3x + (2x - 1) = 9$, so $5x = 10$ and $x = 2$. Then $y = 2(2) - 1 = 3$. Check: $6 + 3 = 9$.
\end{example}

\begin{example}
Solve $2x + 3y = 12$ and $4x - 3y = 6$ by elimination.
\solution The $y$-coefficients are already opposites, so add: $6x = 18$, $x = 3$. Substitute: $6 + 3y = 12$, $y = 2$.
\end{example}

\begin{example}
Where does the line $y = x$ cross $y = -x + 6$?
\solution Set $x = -x + 6$, so $2x = 6$ and $x = 3$; the point is $(3, 3)$. This is exactly the computation a 2D engine does to find where two paths cross.
\end{example}

\section{Exercises}

\begin{multicols}{2}
\begin{exercises}
\item Solve by substitution: $y = x + 4$ and $2x + y = 10$.
\item Solve by elimination: $x + y = 7$ and $x - y = 1$.
\item Solve by elimination: $3x + 2y = 16$ and $2x - 5y = -2$. Show the multipliers you chose.
\item Classify the system $2x - 4y = 8$ and $x - 2y = 4$ as one solution, none, or infinitely many, and explain using slopes.
\item Classify $3x + y = 5$ and $6x + 2y = 7$ the same way.
\item Find the intersection of $y = 2x - 3$ and $y = -x + 6$.
\item Twelve tickets cost \$105. Adult tickets are \$10 and child tickets are \$7. How many of each were sold?
\item You need 20 liters of a 30\% solution by mixing a 10\% solution with a 50\% solution. How many liters of each?
\item For the system in Exercise~5.3, compute $D$ and solve it again with Cramer's rule. Confirm the answers match.
\item Solve by back substitution: $x + y + z = 6$, $2y + z = 7$, $z = 3$.
\end{exercises}
\end{multicols}

\section{Programming Exercises}

\begin{codeexercises}
\item Write \lstinline!solve_2x2(a1, b1, c1, a2, b2, c2)! using Cramer's rule, returning \lstinline!("one", x, y)!, \lstinline!("none",)! or \lstinline!("infinite",)!. When $D$ is zero, decide between none and infinite by checking whether the two equations are proportional. Test with Exercises~5.1 through 5.6.
\item Write \lstinline!intersect(p1, p2, p3, p4)! that takes two lines, each given by two points, and returns their intersection point or \lstinline!None! when they are parallel. Build each line's standard-form coefficients from its two points and reuse C5.1. Verify Exercise~5.6 by picking two points on each line.
\item Write \lstinline!eliminate(rows)! that takes the augmented matrix \lstinline![[a1, b1, c1], [a2, b2, c2]]!, performs elimination with explicit row operations, and prints the matrix after each operation. Run it on Exercise~5.3 and compare the printed steps with your hand work.
\item Write a brute-force checker that scans all integer pairs $(x, y)$ from $-20$ to $20$ and reports every pair satisfying both equations. Use it to confirm C5.1 on every exercise with integer answers, and to show that Exercise~5.4 has many solutions while Exercise~5.5 has none in that range.
\end{codeexercises}
